\chapter{Le articolazioni dell'arto inferiore}

% ---------------------------------------------------------------------------
\section{Generalità e articolazioni del bacino}

Lo scheletro dell'arto inferiore assolve alla funzione di sostenere e
trasmettere a terra il peso del corpo, di permettere la deambulazione
propria della postura bipede e di contribuire al mantenimento
dell'equilibrio. Il \textbf{cingolo pelvico} rappresenta il raccordo con lo
scheletro assile ed è formato dalle ossa dell'ischio, del pube e
dell'ileo; la \textbf{coscia} è costituita dal femore; la \textbf{gamba}
dalla tibia e dalla fibula (o perone); il \textbf{piede} dal tarso, formato
da sette ossa brevi, dal metatarso, formato da cinque ossa lunghe, e dalle
falangi, tre per ciascun dito ad eccezione dell'alluce, che ne possiede
due.

Le articolazioni dell'arto inferiore sono: l'articolazione del bacino,
l'articolazione coxo-femorale, l'articolazione del ginocchio,
l'articolazione tibio-fibulare e le articolazioni del piede.

Il \textbf{bacino} svolge un ruolo strutturale di sostegno, costituisce un
vincolo dinamico fra tronco e arto inferiore, offre ancoraggio ai muscoli
della parete addominale, del dorso e dell'arto inferiore, e funge da
contenitore per gli organi degli apparati digerente, urinario e genitale.
Le articolazioni del bacino sono l'\textbf{articolazione sacroiliaca}, tra
l'osso dell'anca e l'osso sacro, e la \textbf{sinfisi pubica}, tra i due
rami pubici, in rapporto con la linea arcuata. Sull'osso sacro si
riconoscono la base, il promontorio, il processo articolare superiore, la
faccetta auricolare, la tuberosità sacrale, la faccia anteriore e
posteriore e la cresta sacrale laterale; sull'osso dell'anca, a livello
dell'ileo, la tuberosità iliaca, la superficie auricolare, le spine
iliache postero-superiore e postero-inferiore, la grande incisura
ischiatica e il corpo dell'ileo, mentre a livello dell'ischio il corpo
dell'ischio, la spina ischiatica, la piccola incisura ischiatica, il
forame otturatorio, la tuberosità ischiatica e il ramo dell'ischio, in
continuità con il pube.

I legamenti della pelvi comprendono, anteriormente, il legamento
longitudinale anteriore, il legamento ileolombare, i legamenti sacroiliaci
anteriori, il legamento inguinale, il legamento sacrotuberoso e il
legamento sacrospinoso, oltre alla membrana otturatoria; posteriormente, i
legamenti sacroiliaci interossei e i legamenti sacroiliaci posteriori,
tesi tra la cresta iliaca, il tubercolo iliaco, la faccia glutea
dell'ileo, le spine iliache postero-superiore e postero-inferiore e i
processi spinosi delle vertebre lombari, in particolare di L4.

% ---------------------------------------------------------------------------
\section{Articolazione coxo-femorale}

L'articolazione coxo-femorale si forma tra l'acetabolo dell'osso dell'anca
e la testa del femore ed è, come l'articolazione glenomerale della spalla,
un'\textbf{enartrosi} (articolazione a sfera), con tre assi di movimento
perpendicolari fra loro e sei movimenti principali. Sul femore si
riconoscono la testa, il collo anatomico, il grande e il piccolo
trocantere, la linea intertrocanterica e la cresta intertrocanterica,
mentre sull'osso dell'anca si distinguono la cresta iliaca, le spine
iliache antero-superiore, antero-inferiore, postero-superiore e
postero-inferiore, il ciglio cotiloideo dell'acetabolo, il tubercolo
pubico, la spina ischiatica, la tuberosità ischiatica, la tuberosità
glutea e la linea pettinea.

L'articolazione coxo-femorale presenta cinque legamenti principali, che
influiscono sulla stabilità articolare con movimenti di estensione e
rilassamento: il legamento ileofemorale, il legamento ischiofemorale, il
legamento pubofemorale, il legamento rotondo del femore e il legamento
acetabolare trasverso.

Il \textbf{legamento trasverso dell'acetabolo} è una corta banda di
collagene che collega a ponte le due estremità del labbro acetabolare,
chiudendo l'incisura acetabolare e costituendo il margine inferiore della
fossa acetabolare in vivo; nella stessa regione si riconoscono il tetto
dell'acetabolo, la capsula articolare, la faccia semilunare e il legamento
rotondo (della testa) del femore.

I mezzi di unione dell'articolazione coxofemorale comprendono anzitutto la
\textbf{capsula articolare}, che si fissa sul contorno dell'acetabolo e sul
labbro dell'acetabolo e, sul femore, in due punti: anteriormente sulla
linea intertrocanterica e posteriormente sul collo anatomico. I legamenti
di rinforzo dell'articolazione coxo-femorale sono il legamento
ileo-femorale, il legamento pubo-femorale e il legamento
ischio-femorale. Il \textbf{legamento ileo-femorale} origina dal bordo
anteriore dell'osso coxale e si inserisce sul femore lungo la linea
intertrocanterica; prima di raggiungere il femore, il fascio si divide in
due porzioni: il fascio pretrocanterico, il più potente dei legamenti
dell'anca, che raggiunge il margine anteriore del grande trocantere, e il
fascio diretto al piccolo trocantere, che si inserisce nella porzione
inferiore della linea intertrocanterica.

Il \textbf{legamento della testa femorale}, o legamento rotondo, è una
benderella fibrosa molto robusta e appiattita, lunga circa 3-3,5 cm; è
teso dalla fossetta della testa del femore alla fossa dell'acetabolo,
passando al di sotto del legamento trasverso. Al suo interno decorre la
branca posteriore dell'arteria otturatoria; la vascolarizzazione della
testa femorale non avviene tuttavia soltanto tramite il legamento
rotondo, ma anche grazie alle arterie capsulari (tra cui l'arteria
circonflessa mediale e laterale del femore e l'arteria profonda del
femore). Il legamento rotondo è un legamento intrarticolare, avvolto dalla
membrana sinoviale, compreso tra la testa del femore e il batuffolo
adiposo che occupa il fondo dell'acetabolo.

I movimenti dell'articolazione coxo-femorale comprendono la
\textbf{flessione}, che con il ginocchio esteso raggiunge circa 90°, mentre
con il ginocchio flesso può raggiungere i 120° e, in condizioni di massima
mobilità, i 145°; l'\textbf{estensione}, di circa 20°; e
l'\textbf{abduzione}, di circa 30°. A questi movimenti semplici si
associano movimenti combinati, quali l'estensione-adduzione e la
flessione-adduzione della coscia.

% ---------------------------------------------------------------------------
\section{Articolazione del ginocchio}

L'articolazione del ginocchio è un complesso articolare, classificato come
\textbf{ginglimo angolare} (articolazione a troclea), con un asse di
movimento e quindi due movimenti principali. Coinvolge le superfici dei
condili femorali, dei condili tibiali (le cavità glenoidee della tibia) e
della faccia posteriore articolare della patella. Sul femore si
riconoscono l'epicondilo mediale e laterale, i condili mediale e
laterale, la faccia patellare con la relativa gola, la fossa
intercondiloidea e il piano popliteo; sulla tibia la superficie articolare
superiore, i condili mediale e laterale, l'eminenza intercondiloidea, la
tuberosità tibiale e la linea poplitea; sul perone la testa e il collo,
in rapporto con l'articolazione tibiofibulare prossimale; sulla patella la
faccia articolare e l'apice.

\subsection{Menischi e lesioni meniscali}

I \textbf{menischi}, mediale e laterale, aumentano la congruenza fra le
superfici articolari e permettono una distribuzione uniforme delle forze
che gravano sull'articolazione; accompagnano i movimenti dei condili
femorali sulle cavità glenoidee della tibia. Il \textbf{legamento trasverso
del ginocchio} collega tra loro le corna anteriori dei due menischi, per
poi collegarsi alla rotula. Il menisco mediale è connesso posteriormente
anche tramite il legamento meniscofemorale posteriore.

Questi movimenti di adattamento dei menischi sono necessari; a volte,
tuttavia, possono non avvenire correttamente, specie durante rapidi
movimenti di estensione, come nel calcio: i menischi possono allora
essere pizzicati e lesi fra i condili femorali e le cavità glenoidee
tibiali, in particolare il menisco mediale (interno), meno mobile di
quello laterale; si tratta delle cosiddette lesioni meniscali. Dal punto
di vista clinico, la presenza di dolore in rotazione laterale della gamba
indica una lesione al menisco laterale, mentre la presenza di dolore in
rotazione mediale indica una lesione al menisco mediale. Le lesioni
meniscali si distinguono in lesione radiale, lesione a flap, lesione
longitudinale e lesione degenerativa.

Il menisco mediale, fissato più saldamente rispetto a quello laterale, va
incontro a lesioni con maggiore frequenza; tali lesioni sono di solito la
conseguenza di movimenti improvvisi di estensione o rotazione con
l'articolazione del ginocchio in flessione e la gamba bloccata. Tra i
quadri lesionali si riconoscono il menisco a manico di secchio e la
frattura radiale del corno posteriore.

\subsection{Mezzi di unione e borse sinoviali}

La \textbf{capsula fibrosa} riveste tutta l'articolazione del ginocchio. La
\textbf{membrana sinoviale} riveste internamente la membrana fibrosa; a
livello dei menischi si interrompe e si sdoppia, per l'adesione degli
stessi alla membrana fibrosa. Posteriormente la sinovia circonda i
legamenti crociati, che risultano pertanto intracapsulari ma al di fuori
della cavità articolare propriamente detta; sono invece extracapsulari i
legamenti collaterali.

Altre borse sinoviali del ginocchio sono la \textbf{borsa sovrapatellare},
tra il femore e il quadricipite, la \textbf{borsa prepatellare}, tra la
cute e la patella, e la \textbf{borsa infrapatellare}, tra la tibia, il
corpo adiposo infrapatellare e il legamento patellare.

\subsection{Stabilità e legamenti del ginocchio}

L'articolazione del ginocchio è meccanicamente debole e viene
stabilizzata sia dalla potenza e dall'azione dei muscoli e dei tendini che
la circondano (stabilizzatori attivi), sia dai legamenti tesi tra femore e
tibia (stabilizzatori passivi). I legamenti del ginocchio si distinguono
in anteriori, laterali, posteriori e intra-articolari. Tra i legamenti
\textbf{anteriori} rientrano il legamento rotuleo (o patellare) e i
retinacoli laterale e mediale della patella; tra i \textbf{laterali} il
legamento collaterale fibulare e il legamento collaterale tibiale; tra i
\textbf{posteriori} i legamenti poplitei; tra gli
\textbf{intra-articolari} i legamenti crociati anteriore e posteriore.

Il \textbf{legamento patellare} è un robustissimo nastro appiattito che dal
contorno inferiore della rotula si porta alla tuberosità della tibia,
rappresentando il tratto terminale del tendine del muscolo quadricipite
femorale; a livello dell'articolazione femororotulea si realizza
un'\textbf{artrodia}, con scivolamento reciproco fra la troclea femorale e
la faccia posteriore della patella durante i movimenti di
flesso-estensione della femoro-tibiale.

I \textbf{retinacoli laterale e mediale della patella} sono due lamine
fibrose che originano dalle aponeurosi dei muscoli vasto mediale e vasto
laterale (componenti del quadricipite), decorrono ai lati della patella e
si inseriscono ai lati della tuberosità tibiale; sono stabilizzatori della
patella e presentano ispessimenti trasversali e longitudinali.

Il \textbf{legamento collaterale mediale} (o tibiale) è una benderella
fibrosa, aderente alla capsula articolare, che dall'epicondilo mediale del
femore si porta in basso, inserendosi sulla parte più alta della faccia
mediale della tibia. Il \textbf{legamento collaterale laterale} (o
fibulare) è un cordone fibroso che dall'epicondilo laterale del femore si
porta sulla testa della fibula, senza aderire alla capsula articolare. I
legamenti collaterali impediscono i movimenti di lateralità esterna o
interna della tibia (il cosiddetto ginocchio vacillante): il collaterale
mediale impedisce i movimenti di lateralità esterna, cioè l'abduzione
della tibia (movimenti in valgo), mentre il collaterale laterale impedisce
i movimenti di lateralità interna, cioè l'adduzione della tibia
(movimenti in varo). Entrambi sono tesi quando il ginocchio è esteso, in
associazione alla rotazione esterna della tibia.

Il \textbf{legamento popliteo obliquo} è un'espansione distale del muscolo
semimembranoso: si diparte posteriormente dal condilo mediale della tibia
e si porta in alto per inserirsi sul condilo laterale del femore. Il
\textbf{legamento popliteo arcuato} è costituito da due fasci, uno
superiore, diretto al condilo laterale del femore, e uno inferiore, che si
inserisce sulla testa della fibula.

Il \textbf{legamento crociato anteriore} (LCA) origina dalla tibia,
nell'area intercondiloidea anteriore, e si porta in alto e all'indietro
per fissarsi alla faccia mediale del condilo laterale del femore. Il
\textbf{legamento crociato posteriore} è più robusto di quello anteriore:
origina dalla tibia, nell'area intercondiloidea posteriore, e decorre in
senso postero-anteriore, inserendosi sulla faccia laterale del condilo
mediale del femore. I legamenti crociati sono sempre tesi e si detendono
soltanto in parziale flessione del ginocchio e in rotazione esterna della
tibia; la loro funzione principale è impedire i movimenti a cassetto,
cioè lo scorrimento antero-posteriore della tibia rispetto al femore.

Il movimento di \textbf{intrarotazione} della gamba flessa è limitato dai
legamenti crociati, mentre il movimento di \textbf{extrarotazione} è
arrestato dai legamenti collaterali.

\subsection{Nota clinica: rottura dei legamenti crociati}

La rottura del legamento crociato destabilizza l'articolazione del
ginocchio, consentendo alla tibia di muoversi avanti e indietro rispetto
al femore (fenomeno del cassetto, rispettivamente anteriore o
posteriore). La rottura del legamento crociato anteriore si verifica con
un'incidenza circa dieci volte superiore a quella del legamento crociato
posteriore. Il meccanismo più frequente che la determina è il trauma da
intrarotazione forzata con la gamba bloccata; un colpo laterale contro il
ginocchio in estensione completa e il piede bloccato può inoltre causare
la contemporanea rottura del legamento crociato anteriore, del legamento
collaterale mediale e la rottura del menisco mediale, configurando la
cosiddetta ``infelice triade'' di lesione del ginocchio, tipica dei
movimenti di torsione con piede bloccato.

La lesione del legamento crociato anteriore è spesso causata da un colpo
diretto al ginocchio, come nel calcio. La lesione del legamento crociato
posteriore è molto più rara di quella del legamento crociato anteriore ed
è spesso causata da traumi sportivi o da incidenti automobilistici: si
verifica tipicamente in caso di impatto del ginocchio contro il cruscotto
dell'auto durante un incidente stradale (trauma da cruscotto) e in molti
sport di contatto.

\subsection{Cinematica del ginocchio}

L'articolazione del ginocchio consente movimenti di flessione e di
estensione della gamba sulla coscia.

% ---------------------------------------------------------------------------
\section{Articolazione tibio-fibulare}

La \textbf{membrana interossea} della gamba è una lamina fibrosa tesa tra i
margini interossei della tibia e del perone: stabilizza le articolazioni
tibio-fibulari e separa i muscoli anteriori della gamba da quelli
posteriori. Sulla tibia si riconoscono i condili mediale e laterale, la
superficie articolare superiore (cavità glenoidea), l'eminenza
intercondiloidea, la tuberosità tibiale, la linea poplitea, il corpo, le
facce mediale e laterale, il margine anteriore, il solco malleolare e il
malleolo mediale; sul perone la testa, il collo, il corpo e il malleolo
laterale, con la relativa fossa. Distalmente, malleolo mediale e malleolo
laterale formano insieme la pinza malleolare.

L'\textbf{articolazione tibio-fibulare prossimale} è rinforzata dai
legamenti della testa della fibula, anteriori e posteriori.
L'\textbf{articolazione tibio-fibulare distale} è una \textbf{sinartrosi}
tra le estremità distali della tibia e della fibula: la tibia presenta una
faccia incavata a doccia, l'incisura fibulare, che si mette in rapporto
con una superficie rugosa o piana della fibula.

% ---------------------------------------------------------------------------
\section{Le articolazioni del piede}

Le articolazioni del piede comprendono le articolazioni della caviglia, le
articolazioni del tarso, le articolazioni del metatarso e le
articolazioni delle dita.

L'\textbf{articolazione tibio-tarsica} (o talo-crurale) è un
\textbf{ginglimo angolare} (articolazione a troclea), con un asse di
movimento e quindi due movimenti principali: la flessione dorsale
(flessione della caviglia) e la flessione plantare (estensione della
caviglia). È formata dalla pinza malleolare, costituita dal malleolo
mediale e dal malleolo laterale, e dall'astragalo (o talo); a livello del
tarso si riconoscono inoltre il calcagno, l'osso cuboide, il primo, il
secondo e il terzo osso cuneiforme e l'osso scafoide (o osso navicolare),
in rapporto con le ossa metatarsali.

Tra le articolazioni \textbf{intertarsali}, l'astragalo e il calcagno si
articolano anteriormente mediante l'\textbf{articolazione
talo-calcaneo-navicolare} e posteriormente mediante
l'\textbf{articolazione talo-calcaneale}, di tipo trocoide; la
\textbf{calcaneo-cuboidea} è un'articolazione a sella che si stabilisce tra
la faccia anteriore del calcagno e la faccia posteriore del cuboide.

Le \textbf{tarsometatarsali} sono artrodie che uniscono le ossa della fila
distale del tarso con le basi delle ossa metatarsali, consentendo piccoli
movimenti di flessione, estensione e lateralità e permettendo
modificazioni della volta plantare. Le \textbf{intermetatarsali} uniscono
tra loro le basi delle ossa metatarsali, ad esclusione delle prime due,
che sono unite da un legamento interosseo, e consentono piccoli movimenti
di scivolamento. Le \textbf{metatarsofalangee} sono articolazioni a
condilo che uniscono le teste dei metatarsali alle basi delle falangi
prossimali, consentendo movimenti di flessione ed estensione delle dita
del piede. Le \textbf{interfalangee del piede} sono articolazioni a
troclea, due per ciascun dito ad eccezione dell'alluce, che ne possiede
una sola, e uniscono le teste delle falangi alle basi delle stesse,
consentendo movimenti di flessione ed estensione.
